% ============================================
% Section: Discussion
% ============================================
\label{sec:discussion}

\subsection{Design Trade-offs}

The hierarchical hybrid topology is the central design choice and deserves
explicit justification. A purely centralized topology (one fusion server plus
$N$ clients) is simplest and gives a globally optimal map at each step, but it
introduces a single point of failure and a server uplink bottleneck that grows
as $\mathcal{O}(N)$. A purely decentralized P2P topology (full mesh) has no
single point of failure and low latency, but requires $\mathcal{O}(N^2)$ links
and offers no natural place for global consistency. The hierarchical hybrid
captures most of the simplicity of centralization within each cluster while
keeping inter-cluster communication sparse, and it degrades to P2P when the
optional Tier-0 server is absent. The cost is added protocol complexity
(election, heartbeat, submap routing), which we judge acceptable given the
robustness gain.

A second trade-off is the choice of \emph{3DGS as the shared map
representation}. Alternatives include a neural radiance field (NeRF) per node,
or a voxel hash plus a separate Gaussian layer. We chose 3DGS for continuity
with Gaussian-LIC2 (maximizing code reuse and rendering quality), the
existence of mature compression knobs (SH truncation, opacity pruning,
quantization), and the fact that Gaussians are naturally additive and
mergeable---properties a NeRF does not share. The cost is per-Gaussian
bandwidth, which the budget controller manages.

\subsection{Limitations}

\paragraph{Static-scene assumption.} The per-node front-end inherits
Gaussian-LIC2's static-scene assumption. The confidence-weighted merge rule
(Section~\ref{sec:conflict}) is a first-pass treatment of inter-node dynamics
but is not a substitute for proper dynamic-object segmentation across nodes.
A dynamic-aware extension, drawing on RoSe-SLAM \cite{wang2026roseslam} and
Dynamic-LIVO \cite{zhang2026dynamiclivo}, is the most important next step.

\paragraph{No empirical validation yet.} This report presents the design and
analytical bounds only. The bandwidth model assumes ideal compression ratios
and ignores packet-loss retransmission overhead; the robustness claims assume
the failure detectors fire within their nominal timeouts. Both will be tested
in the planned campaign.

\paragraph{Security is minimal.} Signed node descriptors prevent unauthorized
merging, but there is no sybil resistance, no replay protection, and no
confidentiality. A deployment in an untrusted environment would need
additional cryptographic machinery (authenticated routing, per-submap
signatures, encrypted transport).

\paragraph{Heterogeneous clocks.} Online time-offset estimation handles
intra-node clock offsets, but inter-node clock synchronization relies on
overlapping submap registration. Nodes with no spatial overlap for extended
periods may drift in their relative time estimates, degrading merge quality.

\subsection{Open Questions}

Several questions are deferred to future iterations:

\begin{itemize}
  \item \textbf{Dynamic cluster-head handoff.} How to migrate merged state
    when a head moves out of range without stalling the cluster.
  \item \textbf{Cross-node dynamic consensus.} A per-node dynamic-object
    segmenter (RoSe-SLAM-style) plus a cross-node consensus protocol to agree
    on which Gaussians represent moving objects and should be excluded from
    the static map.
  \item \textbf{Energy budget on UAV nodes.} The compressor cost versus
    communication cost trade-off is deployment-specific; UAVs are energy-bound
    and may prefer coarser submaps to save radio power.
  \item \textbf{5G QoS measurement.} A measurement campaign on a real 5G NR
    link to characterize per-flow QoS and tune the budget controller.
  \item \textbf{Sybil resistance.} A decentralized identity mechanism
    (e.g., proof-of-work on node join, or a trusted authority-issued
    certificate) to prevent map poisoning by sybil nodes.
\end{itemize}
